%% ============================================================
%%  Blue Lotus Labs — Methods Section
%%  Suitable for SSRN working paper submission
%%
%%  Compile with:  pdflatex methods.tex  (twice for references)
%%
%%  Required packages: amsmath, amssymb, booktabs, algorithm,
%%                     algorithmicx, algpseudocode, natbib, hyperref
%% ============================================================

\documentclass[12pt]{article}
\usepackage[margin=1.2in]{geometry}
\usepackage{amsmath, amssymb, amsthm}
\usepackage{booktabs}
\usepackage{algorithm}
\usepackage{algpseudocode}
\usepackage{natbib}
\usepackage{hyperref}
\usepackage{xcolor}
\usepackage{setspace}
\onehalfspacing

\newcommand{\E}{\mathbb{E}}
\newcommand{\Q}{\mathbb{Q}}
\newcommand{\P}{\mathbb{P}}
\newcommand{\R}{\mathbb{R}}
\newcommand{\norm}[1]{\left\lVert#1\right\rVert}

\title{\textbf{Constraint-Driven Monte Carlo Stress Testing:\\
A Unified Framework for Institutional Risk Assessment}\\[0.5em]
\large Methods and Model Specification}

\author{Blue Lotus Labs}
\date{\today}

\begin{document}
\maketitle

%% ------------------------------------------------------------------
\section{Overview}
%% ------------------------------------------------------------------

We describe a stress-testing engine that combines three areas of
quantitative risk modelling: hidden Markov regime detection,
extreme value theory (EVT) for tail estimation, and importance
sampling via the Esscher transform.  The engine produces a
distribution of forward-looking drawdown and expected shortfall
(ES) estimates under a well-defined risk-neutral measure, together
with bootstrap confidence intervals and a sensitivity index based
on optimal transport.

%% ------------------------------------------------------------------
\section{Notation}
%% ------------------------------------------------------------------

Let $\{r_t\}_{t=1}^{T}$ be a daily return series, $r_t \in \R$.
Denote by $L_t = -r_t$ the corresponding loss series (positive
losses).  For a random variable $X$ and level $\alpha \in (0,1)$:

\begin{align}
  \text{VaR}_\alpha(X) &= \inf\{x : \P(X \leq x) \geq \alpha\}, \\
  \text{ES}_\alpha(X)  &= \E\!\left[X \mid X \leq \text{VaR}_\alpha(X)\right].
\end{align}

All notation follows \citet{McNeilFreyEmbrechts2015}.

%% ------------------------------------------------------------------
\section{Module 2a — Regime Detection via Hidden Markov Model}
%% ------------------------------------------------------------------

\subsection{Model specification}

We model latent market regimes with a three-state Gaussian Hidden
Markov Model (HMM):
\begin{equation}
  r_t \mid S_t = k \;\sim\; \mathcal{N}(\mu_k,\, \sigma_k^2),
  \quad k \in \{0, 1, 2\},
\end{equation}
where state 0 is \emph{calm} (low volatility), state 1 is
\emph{volatile}, and state 2 is \emph{crisis}.  The hidden state
sequence $\{S_t\}$ follows a first-order Markov chain with
transition matrix $\mathbf{P} \in [0,1]^{3\times 3}$, $\mathbf{P}\mathbf{1}=\mathbf{1}$.

\subsection{Estimation}

Parameters $\Theta = \{\mathbf{P}, \mu_k, \sigma_k^2\}$ are
estimated by maximising the log-likelihood
$\ell(\Theta) = \log p(r_1,\ldots,r_T;\Theta)$
using the Baum--Welch Expectation--Maximisation algorithm
\citep{BaumWelch1972}.  To reduce sensitivity to initialisation,
we run $n_{\text{init}} = 10$ random restarts and retain the
solution with the highest log-likelihood.

State labelling: after estimation, states are ordered by emission
standard deviation $\sigma_k$ (ascending), ensuring consistent
assignment of calm, volatile, and crisis regimes across datasets.

\subsection{Model selection}

Fit quality is assessed via:
\begin{equation}
  \text{AIC} = 2p - 2\ell(\hat\Theta), \qquad
  \text{BIC} = p\log T - 2\ell(\hat\Theta),
\end{equation}
where $p = (K-1) + K(K-1) + 2K$ counts the free parameters
(initial distribution, transition matrix, emission means and
variances) and $K=3$.

The Viterbi algorithm decodes the most probable state sequence
$\hat S_{1:T}$.

%% ------------------------------------------------------------------
\section{Module 2b — Tail Risk via Extreme Value Theory}
%% ------------------------------------------------------------------

\subsection{Peaks-over-Threshold method}

Standard parametric tail models (Student-$t$, normal) underestimate
extreme quantile risk because they lack flexibility in the far tail.
We instead apply the Peaks-over-Threshold (PoT) approach from
Extreme Value Theory \citep{Pickands1975}.

For a threshold $u$ (chosen below), let
$Y_i = L_i - u$ for all $L_i > u$ denote the \emph{exceedances}.
The Pickands--Balkema--de Haan theorem states that, for a
sufficiently high $u$, the distribution of $Y$ converges to a
Generalized Pareto Distribution (GPD):
\begin{equation}
  F_u(y) = 1 - \left(1 + \frac{\xi y}{\beta}\right)^{-1/\xi},
  \quad y > 0,
\end{equation}
with shape $\xi \in \R$ and scale $\beta > 0$.

\subsection{Threshold selection}

We select $u$ using a stability diagnostic on the Hill estimator of
the tail index.  For $k$ order statistics, the Hill estimator is
\citep{Hill1975}:
\begin{equation}
  \hat\xi_k^{\text{Hill}} = \frac{1}{k}\sum_{i=1}^{k}
    \left[\log L_{(i)} - \log L_{(k+1)}\right],
\end{equation}
where $L_{(1)} \geq \cdots \geq L_{(n)}$ are the sorted losses.
The optimal $k^*$ minimises the local variance of $\hat\xi_k^{\text{Hill}}$
over a sliding window, corresponding to the plateau region of the
Hill plot.  We search over $k \in [0.02n,\, 0.15n]$.

\subsection{GPD fitting and analytical risk measures}

Given the optimal threshold $u^* = L_{(k^*+1)}$ and $N_u = k^*$
exceedances, GPD parameters $(\hat\xi, \hat\beta)$ are estimated by
maximum likelihood.

Analytical risk measures follow directly \citep{McNeilFrey2000}:
\begin{align}
  \widehat{\text{VaR}}_\alpha &= u^* + \frac{\hat\beta}{\hat\xi}
    \left[\left(\frac{\alpha n}{N_u}\right)^{-\hat\xi} - 1\right],
    \label{eq:gpd_var}\\
  \widehat{\text{ES}}_\alpha  &= \frac{\widehat{\text{VaR}}_\alpha + \hat\beta - \hat\xi u^*}
    {1 - \hat\xi},
    \label{eq:gpd_es}
\end{align}
for $\hat\xi \neq 0$ and $\hat\xi < 1$ (finite mean condition).
The exponential-tail case ($\hat\xi \to 0$) is handled separately.

%% ------------------------------------------------------------------
\section{Module 2c — Bayesian Shrinkage of Regime Parameters}
%% ------------------------------------------------------------------

Raw within-regime sample means and variances suffer from estimation
error due to the limited number of observations per regime.  We
apply conjugate posterior shrinkage.

\paragraph{Posterior mean.}
Using a Gaussian--Gaussian conjugate model:
\begin{equation}
  \hat\mu_k^* = w_k\,\bar r_k + (1-w_k)\,\mu_0, \qquad
  w_k = \frac{n_k/\sigma^2}{n_k/\sigma^2 + 1/\tau^2},
\end{equation}
where $\mu_0 = \bar r$ is the grand mean, $\sigma^2$ the pooled
variance, and $\tau^2 = \widehat{\mathrm{Var}}(\{\mu_k\})$ the
between-regime variance (ANOVA estimate).

\paragraph{Posterior variance.}
Using the inverse-gamma conjugate update:
\begin{equation}
  \hat\sigma_k^{*2} = \frac{\beta_{\text{IG}} + \frac{1}{2}SS_k}
    {\alpha_{\text{IG}} + \frac{n_k}{2} - 1},
\end{equation}
with hyperparameters $\alpha_{\text{IG}} = 3$, $\beta_{\text{IG}} = \sigma^2(\alpha_{\text{IG}}-1)$.

%% ------------------------------------------------------------------
\section{Change of Measure: Esscher Tilting}
%% ------------------------------------------------------------------

\subsection{Motivation}

Naive Monte Carlo samples paths under the empirical (historical)
measure $\P$.  We instead seek a measure $\Q$ that simultaneously
enforces two constraints on the simulated path distribution:
(i) the mean return equals a target $\mu^*$, and (ii) the expected
shortfall at level $\alpha$ equals the GPD estimate $\text{ES}^*$.

\subsection{Esscher transform}

Following \citet{GerberShiu1994}, we define tilted weights:
\begin{equation}
  w_i \propto \exp\!\left(\lambda_1 \bar r_i
    + \lambda_2 \bar r_i \cdot \mathbf{1}[\bar r_i \leq \widehat{\text{VaR}}_\alpha]\right),
\end{equation}
where $\bar r_i$ is the mean return of path $i$.
The Lagrange multipliers $(\lambda_1, \lambda_2)$ solve the system:
\begin{align}
  \E_\Q[\bar r] &= \mu^*, \\
  \E_\Q[\bar r \mid \bar r \leq \widehat{\text{VaR}}_\alpha] &= \text{ES}^*.
\end{align}

\subsection{Numerical solution}

The system is solved via \texttt{scipy.optimize.fsolve} with five
random starting points.  The solution with the smallest residual
norm is retained; if this norm exceeds $10^{-4}$, the solver falls
back to identity weights ($\lambda_1 = \lambda_2 = 0$), preserving
the simulation under $\P$.

Paths are then importance-resampled with replacement according to
the normalised weights $w_i / \sum_j w_j$, yielding a new path set
distributed under $\Q$.

%% ------------------------------------------------------------------
\section{Module 4 — Stress Metrics and Bootstrap Confidence Intervals}
%% ------------------------------------------------------------------

\subsection{Maximum drawdown}

For each path $\{r_t^{(i)}\}_{t=1}^H$, define the cumulative return
$C_t^{(i)} = \sum_{s=1}^t r_s^{(i)}$ and running maximum
$M_t^{(i)} = \max_{s \leq t} C_s^{(i)}$.  The maximum drawdown is:
\begin{equation}
  D^{(i)} = \min_{1 \leq t \leq H}\left(C_t^{(i)} - M_t^{(i)}\right) \leq 0.
\end{equation}
Key statistics: $\bar D = \frac{1}{N}\sum_i D^{(i)}$,
$D_{0.05} = \text{VaR}_{0.05}(D^{(i)})$, and the 90\% confidence
interval $[D_{0.05}^{\text{dist}}, D_{0.95}^{\text{dist}}]$ across paths.

\subsection{Expected Shortfall on paths}

For each path, per-path ES at level $\alpha$:
\begin{equation}
  \text{ES}_\alpha^{(i)} =
    \frac{1}{\lfloor\alpha H\rfloor}\sum_{t: r_t^{(i)} \leq q_\alpha^{(i)}} r_t^{(i)},
\end{equation}
and the aggregate ES over all paths:
\begin{equation}
  \overline{\text{ES}}_\alpha = \E_\Q\!\left[r_t \mid r_t \leq \text{VaR}_\alpha\right].
\end{equation}

\subsection{Bootstrap confidence intervals}

All scalar metrics $\theta$ (drawdown mean, ES aggregate, recovery time)
are reported with 90\% bootstrap confidence intervals.  We draw
$B = 500$ bootstrap path samples (resampling paths with replacement)
and compute the statistic on each:
\begin{equation}
  \widehat{\text{CI}}_{90\%}(\theta) =
    \left[\hat\theta^*_{(0.05B)},\; \hat\theta^*_{(0.95B)}\right].
\end{equation}

%% ------------------------------------------------------------------
\section{Module 4b — Multi-Asset Extension}
%% ------------------------------------------------------------------

For a portfolio of $N$ assets with weight vector $\mathbf{w}$, we model
the joint return distribution using a multivariate Gaussian:
\begin{equation}
  \mathbf{r}_t \sim \mathcal{N}(\boldsymbol\mu,\, \boldsymbol\Sigma),
\end{equation}
where $\boldsymbol\Sigma$ is estimated via the Ledoit--Wolf shrinkage
estimator \citep{LedoitWolf2004}:
\begin{equation}
  \hat{\boldsymbol\Sigma}_{\text{LW}} =
    (1-\rho)\,\mathbf{S} + \rho\,\mathbf{F},
\end{equation}
with $\mathbf{S}$ the sample covariance, $\mathbf{F}$ a structured target
(scaled identity), and $\rho^*$ the analytic shrinkage intensity.

Correlated paths are generated via Cholesky decomposition:
$\mathbf{L} = \text{chol}(\hat{\boldsymbol\Sigma}_{\text{LW}})$,
$\mathbf{r}_t^{(i)} = \boldsymbol\mu + \mathbf{L}\mathbf{z}_t^{(i)}$,
$\mathbf{z}_t^{(i)} \sim \mathcal{N}(\mathbf{0},\mathbf{I})$.

Portfolio returns are $p_t^{(i)} = \mathbf{w}^\top \mathbf{r}_t^{(i)}$,
and all stress metrics are computed on $\{p_t^{(i)}\}$.

%% ------------------------------------------------------------------
\section{Module 5 — Model Fragility Index}
%% ------------------------------------------------------------------

\subsection{Definition}

The Model Fragility Index (MFI) quantifies the sensitivity of the
risk estimates to small parameter perturbations using the
Wasserstein-1 (Earth Mover's) distance \citep{Villani2008}.

Let $D_0 = \{D^{(i)}\}$ be the baseline drawdown distribution and
$D_\theta$ the distribution under perturbation $\theta \in \Theta$.
Define:
\begin{equation}
  \text{MFI} = \frac{1}{|\Theta|}\sum_{\theta\in\Theta}
    \frac{W_1(D_0,\, D_\theta)}{|\overline{\text{ES}}_0|},
\end{equation}
where
$W_1(\mu,\nu) = \int_0^1 |F_\mu^{-1}(u) - F_\nu^{-1}(u)|\,du$
is the Wasserstein-1 distance between the two CDFs and
$|\overline{\text{ES}}_0|$ normalises for scale.

\subsection{Perturbation set $\Theta$}

\begin{center}
\begin{tabular}{lll}
\toprule
Name & Parameter & Perturbation \\
\midrule
\texttt{mean\_+10pct}  & Regime means    & $\times 1.10$ \\
\texttt{mean\_-10pct}  & Regime means    & $\times 0.90$ \\
\texttt{vol\_+20pct}   & Regime std devs & $\times 1.20$ \\
\texttt{vol\_-20pct}   & Regime std devs & $\times 0.80$ \\
\texttt{trans\_±15pct} & Transition matrix & Dirichlet noise $\pm 15\%$ \\
\bottomrule
\end{tabular}
\end{center}

\subsection{Interpretation}

\begin{center}
\begin{tabular}{lll}
\toprule
MFI range & Grade & Interpretation \\
\midrule
$[0,\, 0.25)$ & Robust   & Estimates stable across perturbations \\
$[0.25, 0.55)$ & Moderate & Some sensitivity to model inputs \\
$[0.55, \infty)$ & Fragile  & Estimates sensitive; interpret with caution \\
\bottomrule
\end{tabular}
\end{center}

This formalises the intuition from \citet{ContDeguest2013} that a
``good'' risk model should not change its conclusions under small,
plausible misspecifications of its inputs.

%% ------------------------------------------------------------------
\section{Module 5b — Historical Backtest Validation}
%% ------------------------------------------------------------------

We validate the engine's predictions against three documented stress
episodes (Table~\ref{tab:backtest_periods}).

\begin{table}[h]
\centering
\caption{Stress periods used for backtest validation.}
\label{tab:backtest_periods}
\begin{tabular}{lll}
\toprule
Period & Start & End \\
\midrule
2008 Financial Crisis  & Sep 2007 & Mar 2009 \\
March 2020 COVID Crash & Feb 2020 & Apr 2020 \\
2022 Rate Shock        & Jan 2022 & Oct 2022 \\
\bottomrule
\end{tabular}
\end{table}

For each period $p$ with $n_p$ trading days, we compute:
\begin{itemize}
  \item \textbf{Realized max drawdown}: $D_p^{\text{real}} = \min_t(C_t^p - M_t^p)$.
  \item \textbf{Predicted} 5th-percentile drawdown $D_{0.05}$ from the
    full-sample engine run.
  \item \textbf{Coverage}: whether $D_p^{\text{real}} \in
    [D_{0.05}^{\text{pred}}, D_{0.95}^{\text{pred}}]$.
  \item \textbf{Coverage ratio}: $D_p^{\text{real}} / D_{0.05}^{\text{pred}}$
    (ratio $> 1$ indicates under-estimation of risk; $< 1$ indicates
    over-estimation).
\end{itemize}

\noindent\textit{Note:} This is an in-sample backtest since the engine
is fitted on the full return history which includes the stress periods.
A rigorous out-of-sample evaluation would re-fit the model on data up
to the start of each period; this is left for future work.

%% ------------------------------------------------------------------
\section{Summary of Computational Complexity}
%% ------------------------------------------------------------------

\begin{center}
\begin{tabular}{ll}
\toprule
Step & Complexity \\
\midrule
HMM fitting (Baum-Welch, $n_{\text{init}}$ restarts) & $O(n_{\text{init}}\cdot T \cdot K^2)$ \\
GPD threshold search & $O(T \log T)$ \\
Monte Carlo generation ($N$ paths, $H$ steps) & $O(NH)$ \\
Esscher solver (5 starting points) & $O(N)$ \\
Bootstrap CIs ($B$ resamples) & $O(BNH)$ \\
Wasserstein MFI (5 perturbations) & $O(5 \cdot NH + 5 \cdot N \log N)$ \\
\bottomrule
\end{tabular}
\end{center}

For $T=3{,}800$ (SPY 2010--2025), $N=10{,}000$, $H=252$, $B=500$,
the full pipeline completes in approximately 60--90 seconds on a
single CPU core.

%% ------------------------------------------------------------------
\section{Limitations and Extensions}
%% ------------------------------------------------------------------

Several limitations warrant attention.

\paragraph{Stationarity.}  The HMM assumes stationary transition
probabilities and emission distributions.  In practice, these
parameters shift over time (e.g., post-2008 regulatory changes alter
the probability of crisis regimes).  Time-varying HMMs or
regime-switching GARCH models could address this.

\paragraph{Gaussian emissions.}  Within each regime, returns are
assumed Gaussian.  While the GPD tail corrects for this at the extreme
quantile level, the body of the regime distributions remains
symmetrically parametrised.  Asymmetric emission models (e.g.,
skew-$t$) would improve fit at intermediate quantiles.

\paragraph{Single-factor regimes.}  The HMM detects regimes based on
univariate return dynamics.  A multivariate HMM incorporating volume,
implied volatility, or credit spreads as additional signals could
improve regime detection quality.

\paragraph{In-sample backtest.}  The historical backtest in Module 5b
uses the full return history for fitting.  Walk-forward validation
would provide a stronger empirical test of predictive performance.

%% ------------------------------------------------------------------
\bibliographystyle{plainnat}
\begin{thebibliography}{99}

\bibitem[Baum \& Welch(1972)]{BaumWelch1972}
Baum, L.~E. and Welch, G.~E. (1972).
\newblock An inequality and associated maximization technique in statistical
  estimation for probabilistic functions of Markov processes.
\newblock \emph{Inequalities}, 3, 1--8.

\bibitem[Cont \& Deguest(2013)]{ContDeguest2013}
Cont, R. and Deguest, R. (2013).
\newblock Robustness and sensitivity analysis of risk measurement procedures.
\newblock \emph{Quantitative Finance}, 10(6), 593--606.

\bibitem[Gerber \& Shiu(1994)]{GerberShiu1994}
Gerber, H.~U. and Shiu, E.~S.~W. (1994).
\newblock Option pricing by Esscher transforms.
\newblock \emph{Transactions of the Society of Actuaries}, 46, 99--191.

\bibitem[Hill(1975)]{Hill1975}
Hill, B.~M. (1975).
\newblock A simple general approach to inference about the tail of a
  distribution.
\newblock \emph{The Annals of Statistics}, 3(5), 1163--1174.

\bibitem[Ledoit \& Wolf(2004)]{LedoitWolf2004}
Ledoit, O. and Wolf, M. (2004).
\newblock A well-conditioned estimator for large-dimensional covariance
  matrices.
\newblock \emph{Journal of Multivariate Analysis}, 88(2), 365--411.

\bibitem[McNeil \& Frey(2000)]{McNeilFrey2000}
McNeil, A.~J. and Frey, R. (2000).
\newblock Estimation of tail-related risk measures for heteroscedastic
  financial time series: an extreme value approach.
\newblock \emph{Journal of Empirical Finance}, 7(3-4), 271--300.

\bibitem[McNeil, Frey \& Embrechts(2015)]{McNeilFreyEmbrechts2015}
McNeil, A.~J., Frey, R. and Embrechts, P. (2015).
\newblock \emph{Quantitative Risk Management: Concepts, Techniques and Tools},
  revised edition.
\newblock Princeton University Press.

\bibitem[Pickands(1975)]{Pickands1975}
Pickands, J. (1975).
\newblock Statistical inference using extreme order statistics.
\newblock \emph{The Annals of Statistics}, 3(1), 119--131.

\bibitem[Villani(2008)]{Villani2008}
Villani, C. (2008).
\newblock \emph{Optimal Transport: Old and New}.
\newblock Springer.

\end{thebibliography}

\end{document}
